%% Rechnungsvorlage „prag“ (ADR-0011): klassisch-tabellarisch und schwarz-weiß, die Ordnung trägt das Gitter.
\documentclass[10pt,a4paper]{article}

%% Als Erstes, denn die Grundlage lädt pdfx; was sie lädt, lädt diese Vorlage nicht noch einmal (ADR-0022).
\usepackage{grundlage}

\usepackage{geometry}
%% xltabular lädt tabularx selbst, das für die Absetzung bleibt.
\usepackage{xltabular}
\usepackage{fancyhdr}
\usepackage{array}
\usepackage{etoolbox}  %% \ifblank für die Angaben, die leer sein dürfen
\usepackage{tikz}      %% das Logo

%% Gruppiert ab vier Stellen; Trennzeichen und Stellung des Währungszeichens setzt die Sprachdatei (ADR-0021).
\sisetup{group-minimum-digits = 4}

\newcommand{\waehrung}{<<WAEHRUNG>>}

%% Nach etoolbox und siunitx, denn die Sprachdatei braucht beide.
\input{sprache}

%% Fira Sans, aus TeX Live.
\setmainfont{FiraSans}[
  Extension      = .otf,
  UprightFont    = *-Regular,
  BoldFont       = *-Bold,
  ItalicFont     = *-Italic,
  BoldItalicFont = *-BoldItalic,
]

%% includefoot misst den unteren Rand bis zur Grundlinie des Fußes, und die ist seine Unterkante (siehe \fuss).
%% 16mm-\maxdepth hält sie, wo sie ohne Seitenfuß stand. \footskip setzt \fuss selbst, unten.
\geometry{a4paper, left=20mm, right=20mm, top=18mm, bottom=\dimexpr 16mm-\maxdepth\relax,
  includefoot, headheight=12pt, headsep=6mm}

\setlength{\parindent}{0pt}
\fancypagestyle{erste}{\fancyhf{}\renewcommand{\headrulewidth}{0pt}\fancyfoot[C]{\fuss}\fancyfoot[R]{\fortsetzung}}
\fancypagestyle{folge}{\fancyhf{}\renewcommand{\headrulewidth}{0pt}%
  \fancyhead[L]{\small <<DOKUMENTART>> <<RECHNUNGSNUMMER>> · <<RECHNUNGSDATUM>>}%
  \fancyhead[R]{\small \wortseitevon{\thepage}{\pageref*{LastPage}}}%
  \fancyfoot[C]{\fuss}\fancyfoot[R]{\fortsetzung}}
\pagestyle{folge}

%% Im Fuß, damit er an jeder Seite hängt, ohne Raum einzunehmen: 3 mm unter dem Ende des Textbereichs.
\newcommand{\fortsetzung}{\wennweitereseite{%
  \raisebox{\dimexpr\footskip-3mm\relax}[0pt][0pt]{\small \wortfortsetzung}}}
\raggedright

%% Keine Silbentrennung: Eigennamen, Beträge und Anschriften.
\hyphenpenalty=10000
\exhyphenpenalty=10000
\tolerance=2000

%% Etwas Luft in den Zeilen des Gitters, sonst kleben die Ziffern an den Linien.
\renewcommand{\arraystretch}{1.35}

%% Ein Code ist breiter als €, um so viel wachsen die drei Betragsspalten: So hat ein Betrag in jeder Währung Platz
%% wie in Euro (#480). Mit \normalfont, denn vor \begin{document} gilt die Schrift des Designs noch nicht.
\newlength{\zeichenmehrbreite}\settowidth{\zeichenmehrbreite}{\normalfont\betrag{0}}
\newlength{\eurozeichenbreite}\settowidth{\eurozeichenbreite}{\normalfont\eurobetrag{0}}
\addtolength{\zeichenmehrbreite}{-\eurozeichenbreite}
%% $ und £ sind schmaler als €, trotzdem schrumpft keine Spalte: Sonst sähe jede Rechnung in US-Dollar anders aus.
\ifdim\zeichenmehrbreite<0pt \setlength{\zeichenmehrbreite}{0pt}\fi

%% Dieselbe Breite in allen drei Tabellen, sonst fluchten ihre Linien nicht. Nachgemessen braucht „−19.040,00 €“
%% rund 21 mm.
\newlength{\betragsbreite}\setlength{\betragsbreite}{\dimexpr 24mm+\zeichenmehrbreite\relax}
\newcolumntype{B}{>{\raggedleft\arraybackslash}p{\betragsbreite}}

%% Die Summen stehen ohne Linie untereinander, erst die Abschlusszeile setzt wieder eine.
\newcommand{\positionstrenner}{\\\hline}
\newcommand{\summentrenner}{\\}
\newcommand{\zeilenwechsel}{\\}
\newcommand{\zusammen}[1]{\mbox{#1}}
%% Das Land unter der Anschrift, in Versalien nach DIN 5008.
\newcommand{\landzeile}[1]{\MakeUppercase{#1}}
%% \betrag und \prozent stehen in der Sprachdatei.
\newcommand{\menge}[1]{\num{#1}}

%% Ist der Wert leer, entfällt die Zeile samt Beschriftung.
\newcommand{\kopfangabe}[2]{\ifblank{#2}{}{#1 & #2\\}}

%% Nur für Sätze des Programms, keine Eingabedaten: \programmsatz setzt den breiten Satzabstand.
\newcommand{\hinweisabsatz}[1]{\par\vspace{5mm}\noindent%
  \ungeteilt{\programmsatz{#1}}}

%% Ohne \programmsatz: Der Befreiungsgrund kommt aus den Eingabedaten (#143).
\newcommand{\steuerhinweis}[1]{\noindent\ungeteilt{#1}}

\newcommand{\absetzungskopf}[1]{\par\vspace{6mm}\noindent #1\par\vspace{1.5mm}}

%% Der Renderer liefert die Zeilen ohne Zeilenende nach der letzten, das setzt diese Tabelle.
\newcommand{\absetzungstabelle}[1]{\par\noindent%
  \begin{tabularx}{\linewidth}{|>{\raggedright\arraybackslash}X|B|}\hline#1\\\hline\end{tabularx}\par}

%% \newline statt \\, das in der X-Spalte die Tabellenzeile beendete. \normalfont, weil die Bezeichnung einer
%% Position fett steht.
\newcommand{\positionsdetail}[1]{\newline{\normalfont\small #1}}
\newcommand{\absetzungsdetail}[1]{\newline{\small #1}}

%% Einmal gefragt, danach nur der Schalter: Zeile, Spaltentyp und \spaltenzahl müssen dieselbe Antwort
%% bekommen, sonst hat eine Zeile eine Zelle mehr, als die Tabelle Spalten hat.
\newtoggle{rabattspalte}
\ifblank{<<RABATTSPALTE>>}{}{\toggletrue{rabattspalte}}
\newcommand{\mitrabatt}[1]{\iftoggle{rabattspalte}{#1}{}}
\iftoggle{rabattspalte}{\newcolumntype{R}{>{\raggedleft\arraybackslash}p{\dimexpr 14mm+\zeichenmehrbreite\relax}|}}{\newcolumntype{R}{}}

%% \bfseries statt \textbf: Eine Leerzeile in der Beschreibung ist ein Absatz, und den nimmt \textbf nicht an (#488).
\newcommand{\positionszeile}[7]{#1 & {\bfseries #2} & #3 & #4 & #5\mitrabatt{ & #6} & #7}

\newcommand{\positionskopf}{\hline
  \textbf{\wortpos} & \textbf{\wortbezeichnung} & \textbf{\wortmenge} & \textbf{\worteinheit} & \textbf{\worteinzelpreis}\mitrabatt{ & \textbf{\wortrabatt}} & \textbf{\wortgesamt}\\
  \hline}

%% Muss im ersten Argument von \multicolumn zur Zahl expandieren; \iftoggle tut das.
\newcommand{\spaltenzahl}{\iftoggle{rabattspalte}{7}{6}}

\newcommand{\summenzeile}[2]{#1 & #2}

%% Beendet die Zeile davor selbst (#147).
\newcommand{\abschlusszeile}[2]{\\\hline\textbf{#1} & \textbf{#2}}

%% prag nennt die Steuer ausgeschrieben und den Satz dahinter.
\newcommand{\nettozeile}[1]{\summenzeile{\wortzwischensumme}{#1}}
\newcommand{\steuerzeile}[3]{\summenzeile{\wortsteuer{#1}{#2}}{#3}}
\newcommand{\gesamtzeile}[1]{\abschlusszeile{\wortgesamtbetrag}{#1}}

%% Die Bankverbindung steht im Fuß, als Angabe des Briefbogens und nicht als Zahlungsaufforderung. \bankblock
%% bleibt deshalb leer (ADR-0011).
\newcommand{\bankblock}{}

%% Das Logo definiert \logozeichnung und optional \logoskalierung (ADR-0012). TikZ und keine Bilddatei:
%% luatex.def gibt jedem per \includegraphics eingebundenen PDF eine Transparenzgruppe mit, und die verbietet
%% PDF/A-1 (#197).
\newlength{\logohoehe}\setlength{\logohoehe}{12mm}
\newlength{\versalhoehe}
\ifblank{<<LOGO>>}{}{\input{<<LOGO>>}}
\providecommand{\logoskalierung}{1}

%% ------------------------------------------------------------------ Fuß
%% \raisebox{\depth} macht die Unterkante zur Grundlinie: So wächst der Fuß nach oben und nie in den Rand.
\newcommand{\fuss}{\raisebox{\depth}{\parbox[t]{\textwidth}{\footnotesize
\noindent
\begin{minipage}[t]{0.32\linewidth}\raggedright
  <<FIRMA>>\zeilenwechsel{}<<STRASSE>>\zeilenwechsel{}<<PLZ>> <<ORT>>\ifblank{<<WEBSEITE>>}{}{\zeilenwechsel{}<<WEBSEITE>>}
\end{minipage}%
\hfill
\begin{minipage}[t]{0.34\linewidth}\raggedright
  \wortustid: <<USTID>>\ifblank{<<STEUERNUMMER>>}{}{\\\wortsteuernummer: <<STEUERNUMMER>>}\ifblank{<<REGISTERGERICHT>>}{}{\\\wortregistergericht: <<REGISTERGERICHT>>\ifblank{<<REGISTERNUMMER>>}{}{, <<REGISTERNUMMER>>}}\ifblank{<<GESCHAEFTSFUEHRUNG>>}{}{\\\wortgeschaeftsfuehrung: <<GESCHAEFTSFUEHRUNG>>}
\end{minipage}%
\hfill
\begin{minipage}[t]{0.28\linewidth}\raggedright
  <<INSTITUT>>\\IBAN: <<IBAN>>\\BIC: <<BIC>>
\end{minipage}}}}

%% \footskip aus der Höhe des Fußes: Ein fester Wert wäre bei einem umbrechenden Firmennamen zu klein, und
%% fancyhdr schöbe den Fuß nach unten.
\newlength{\fusshoehe}\settoheight{\fusshoehe}{\fuss}
\geometry{footskip=\dimexpr\fusshoehe+\footruleskip+4mm\relax}

\begin{document}
\thispagestyle{erste}

%% ------------------------------------------------------- Titel und Kopfangaben
%% Das Logo nimmt keinen Raum ein ([0pt][0pt]), so bleibt alles, wo es ohne Logo stand. Höher als der Block der
%% Kopfangaben ragte es in die Anschriften; mit \logoskalierung liegt die Grenze bei gut 2.
\settoheight{\versalhoehe}{\LARGE\bfseries M}%
\noindent\ifblank{<<LOGO>>}{}{\raisebox{\dimexpr\versalhoehe-\height\relax}[0pt][0pt]{%
  \resizebox{!}{\logoskalierung\logohoehe}{\begin{tikzpicture}\logozeichnung\end{tikzpicture}}}}%
\hfill
\begin{tabular}[t]{@{}r@{\hspace{3mm}}l@{}}
  \multicolumn{2}{@{}r@{}}{\LARGE\bfseries <<DOKUMENTART>>}\\[2mm]
  \kopfangabe{\wortrechnungsnr}{<<RECHNUNGSNUMMER>>}
  \kopfangabe{\wortkundennr}{<<KUNDENNUMMER>>}
  \kopfangabe{\wortauftragsnr}{<<AUFTRAGSNUMMER>>}
  \kopfangabe{\wortdatum}{<<RECHNUNGSDATUM>>}
  \kopfangabe{<<ZEITRAUM_BESCHRIFTUNG>>:}{<<LEISTUNGSZEITRAUM>>}
  \kopfangabe{\wortvereinnahmtam}{<<VEREINNAHMT_AM>>}
  \kopfangabe{\wortreferenz}{<<KUNDENREFERENZ>>}
\end{tabular}

\vspace{12mm}

%% ------------------------------------------------------ Empfänger und Absender
%% Die Person unter der Firmierung nur ohne Geschäftsführung, die hier im Fuß steht.
\begin{minipage}[t]{0.5\linewidth}\raggedright
  <<EMPFAENGER>>\ifblank{<<KUNDE_USTID>>}{}{\\[2mm]\wortustid: <<KUNDE_USTID>>}
\end{minipage}%
\hfill
\begin{minipage}[t]{0.45\linewidth}\raggedleft
  <<FIRMA>>\ifblank{<<GESCHAEFTSFUEHRUNG>>}{\zeilenwechsel{}<<INHABER>>}{}\zeilenwechsel{}<<STRASSE>>\zeilenwechsel{}<<PLZ>> <<ORT>>\ifblank{<<TELEFON>>}{}{\zeilenwechsel{}\worttelefon: <<TELEFON>>}\zeilenwechsel{}<<EMAIL>>
\end{minipage}

\vspace{12mm}

%% --------------------------------------------------------------- Einleitung
%% Aus den Eingabedaten, deshalb nicht in \hinweisabsatz (#143).
\ifblank{<<EINLEITUNG>>}{}{\noindent <<EINLEITUNG>>\par\vspace{5mm}}

%% --------------------------------------------------------- Positionen, Summen
%% Jede \hline ist in longtable eine Stelle, an der die Seite brechen darf (ADR-0020). Die Summen sind die letzte
%% Zeile, damit sie nie ohne Position darüber stehen; die Linie davor ist deshalb ein \hrule und keine \hline,
%% und \\* änderte daran nichts. \LTpre und \LTpost halten den Abstand der früheren tabularx.
\setlength{\LTpre}{\lineskip}\setlength{\LTpost}{0pt}%
\setlength{\LTleft}{0pt}\setlength{\LTright}{0pt}%
\begin{xltabular}{\linewidth}{|>{\centering\arraybackslash}p{7mm}|>{\raggedright\arraybackslash}X|>{\raggedleft\arraybackslash}p{13mm}|>{\raggedright\arraybackslash}p{15mm}|>{\raggedleft\arraybackslash}p{\dimexpr 22mm+\zeichenmehrbreite\relax}|RB|}
  \positionskopf
  \endfirsthead
  \positionskopf
  \endhead
<<POSITIONEN>>\\*
  \noalign{\hrule height\arrayrulewidth}
  \multicolumn{\spaltenzahl}{@{}l@{}}{\begin{tabularx}{\linewidth}[t]{|>{\raggedright\arraybackslash}X|B|}
<<SUMMEN>>\\
  \hline
  \end{tabularx}}
\end{xltabular}
\nointerlineskip\vspace{\lineskip}

%% ------------------------------------------------------ Die fachlichen Sätze
%% Die Zahlungsfrist zwischen Absetzung und Skonto: erst was offen ist, dann wann, dann der mögliche Abzug (#211).
<<ARTHINWEIS>>

%% Den Abstand der Steuerhinweise setzt die Vorlage, und nur, wenn etwas dasteht (#220).
\ifblank{<<STEUERHINWEISE>>}{}{\par\vspace{5mm}<<STEUERHINWEISE>>}

<<UMRECHNUNG>><<ABSETZUNG>><<ZAHLUNGSFRIST>><<SKONTO>>

<<BANKVERBINDUNG>>

%% Der Schlusssatz hängt am Absatz davor, sonst stünde er allein auf der letzten Seite. \nopagebreak vor dem
%% Abstand: Hinter einer Strafe ist der Abstand keine Bruchstelle mehr.
\nopagebreak
\vspace{8mm}

\ungeteilt{<<SCHLUSSSATZ>>}

\end{document}
